\documentclass[lualatex,ja=standard,a4paper,12pt]{bxjsarticle}
\usepackage{amsmath, amssymb, amsthm}
\usepackage{geometry}
\usepackage{xcolor}
\usepackage{tcolorbox}
\usepackage{booktabs}
\usepackage{array}
\usepackage{pgfplots}
\pgfplotsset{compat=1.18}

\geometry{margin=25mm}

\tcbuselibrary{theorems, skins, breakable}

\newtcolorbox{defbox}[1]{
  colback=blue!5, colframe=blue!60!black,
  fonttitle=\bfseries, title={#1}, arc=3pt, breakable
}
\newtcolorbox{thmbox}[1]{
  colback=green!5, colframe=green!50!black,
  fonttitle=\bfseries, title={#1}, arc=3pt, breakable
}
\newtcolorbox{pointbox}{
  colback=orange!8, colframe=orange!70!black,
  fonttitle=\bfseries, title={ポイント}, arc=3pt, breakable
}
\newtcolorbox{exbox}[1]{
  colback=gray!6, colframe=gray!50,
  fonttitle=\bfseries, title={例：#1}, arc=3pt, breakable
}
\newtcolorbox{derivbox}{
  colback=purple!4, colframe=purple!50!black,
  fonttitle=\bfseries, title={導出}, arc=3pt, breakable
}

\title{{\large Hoff 第3章} \\[6pt]
       \textbf{3.3\; Exponential Families and Conjugate Priors} \\[4pt]
       {\normalsize ゼミ資料}}
\author{}
\date{}

\begin{document}
\maketitle
\thispagestyle{empty}

% ===========================
\section{この節のモチベーション}
% ===========================

第3章ではここまで，二項モデルとポアソンモデルをそれぞれ独立に扱い，
それぞれについて共役事前分布（ベータ分布・ガンマ分布）を構成してきた．

\medskip
\noindent\textbf{疑問：} これらは「たまたま」上手くいったのか，
それとも共通した構造があるのか？

\medskip
\noindent 答えは後者で，両モデルは\textbf{1次元指数型分布族}という
共通のクラスに属している．この構造を使えば，
共役事前分布を一般的に構成できる．

% ===========================
\section{1次元指数型分布族}
% ===========================

\begin{defbox}{1次元指数型分布族}
密度（または確率質量関数）が次の形で書けるモデルを
\textbf{1次元指数型分布族}と呼ぶ：
\[
  p(y|\phi) = h(y)\,c(\phi)\,e^{\phi\, t(y)},
\]
ここで
\begin{itemize}
  \item $\phi$：自然パラメータ（推定したいパラメータ）
  \item $t(y)$：十分統計量
  \item $c(\phi)$：正規化定数（$\phi$ のみに依存）
  \item $h(y)$：$\phi$ に依存しない部分
\end{itemize}
\end{defbox}

\begin{pointbox}
指数型分布族の最大の特徴は「データの情報が $t(y)$ に集約される」こと．
$n$ 個の i.i.d. サンプル $y_1,\ldots,y_n$ に対しては，
十分統計量 $\sum t(y_i)$（または $\bar{t}(y) = \frac{1}{n}\sum t(y_i)$）
だけ知っていれば $\phi$ についての推論に十分である．
\end{pointbox}

% ===========================
\section{共役事前分布の一般形}
% ===========================

Diaconis and Ylvisaker (1979) は，指数型分布族に対する共役事前分布が
次の形で一般的に書けることを示した：

\begin{thmbox}{共役事前分布（一般形）}
\[
  p(\phi \mid n_0, t_0) = \kappa(n_0, t_0)\,c(\phi)^{n_0}\,e^{n_0 t_0 \phi},
\]
ここで $n_0 > 0$，$t_0$ はハイパーパラメータ．
\end{thmbox}

この式は一見複雑に見えるが，
$p(y|\phi) = h(y)c(\phi)e^{\phi t(y)}$ と見比べると，
「$n_0$ 個のサンプルの尤度」と同じ形をしていることがわかる．

% ===========================
\section{事後分布の導出}
% ===========================

$Y_1,\ldots,Y_n \sim$ i.i.d. $p(y|\phi)$ のもとで，
上の共役事前分布とベイズの定理を使うと：

\begin{derivbox}
\begin{align*}
  p(\phi \mid y_1,\ldots,y_n)
  &\propto p(\phi) \times p(y_1,\ldots,y_n|\phi) \\
  &\propto c(\phi)^{n_0} e^{n_0 t_0 \phi}
   \times \prod_{i=1}^n c(\phi)\,e^{\phi\,t(y_i)} \\
  &= c(\phi)^{n_0+n}
     \exp\!\left\{\phi \times \left(n_0 t_0 + \sum_{i=1}^n t(y_i)\right)\right\} \\
  &\propto p\!\left(\phi \;\middle|\; n_0 + n,\;\;
    \frac{n_0 t_0 + n\bar{t}(y)}{n_0 + n}\right).
\end{align*}
ここで $\bar{t}(y) = \frac{1}{n}\sum_{i=1}^n t(y_i)$．
\end{derivbox}

\begin{thmbox}{事後分布}
\[
  \boxed{
    p(\phi \mid y_1,\ldots,y_n)
    = p\!\left(\phi \;\middle|\; n_0 + n,\;\;
      \frac{n_0 t_0 + n\bar{t}(y)}{n_0 + n}\right)
  }
\]
事前分布と同じクラスの分布になる（共役性の確認）．
\end{thmbox}

事後分布のパラメータを読むと，更新則が明快に見える：

\medskip
\begin{center}
\renewcommand{\arraystretch}{1.6}
\begin{tabular}{lll}
\toprule
& 事前 & 事後 \\
\midrule
サンプルサイズ & $n_0$ & $n_0 + n$ \\
十分統計量の平均 & $t_0$ & $\dfrac{n_0 t_0 + n\bar{t}(y)}{n_0 + n}$ \\
\bottomrule
\end{tabular}
\end{center}

\medskip
事後の十分統計量平均は，\textbf{事前の $t_0$ とデータの $\bar{t}(y)$ の加重平均}になっている．

% ===========================
\section{ハイパーパラメータの解釈}
% ===========================

\subsection*{$n_0$ の解釈（事前サンプルサイズ）}

事前分布 $p(\phi|n_0, t_0)$ は，$\phi$ の関数として見ると
「$t(y_i)/n_0 = t_0$ となる $n_0$ 個の仮想的な事前観測から得られた
尤度」と同じ形をしている．

つまり $n_0$ は「事前情報の強さ」を，
データ数に換算して表したものである．

\subsection*{$t_0$ の解釈（$t(Y)$ の事前期待値）}

Diaconis and Ylvisaker (1979) は次を示した：

\[
  E[t(Y)] = E\!\bigl[E[t(Y)|\phi]\bigr]
  = E\!\left[-\frac{c'(\phi)}{c(\phi)}\right] = t_0.
\]

すなわち，$t_0$ は十分統計量 $t(Y)$ の\textbf{事前期待値}を表す．

\begin{pointbox}
両パラメータの直感的な解釈：
\begin{itemize}
  \item $n_0$：「事前観測が何個分の情報量か」
  \item $t_0$：「十分統計量 $t(Y)$ はどのくらいの値になると思っているか」
\end{itemize}
これは5.5節で正規モデルに拡張される発想の原型である．
\end{pointbox}

% ===========================
\section{例：二項モデル}
% ===========================

\begin{exbox}{二項モデル}

$Y \sim \mathrm{Bernoulli}(\theta)$ の密度を指数型分布族の形に整理する：

\begin{align*}
  p(y|\theta)
  &= \theta^y(1-\theta)^{1-y} \\
  &= \left(\frac{\theta}{1-\theta}\right)^y (1-\theta) \\
  &= e^{\phi y}(1+e^\phi)^{-1},
\end{align*}

ここで $\phi = \log[\theta/(1-\theta)]$ は\textbf{対数オッズ（log-odds）}．

\medskip
\noindent\textbf{各成分の対応：}

\medskip
\begin{center}
\begin{tabular}{ll}
\toprule
記号 & 内容 \\
\midrule
$\phi$ & $\log[\theta/(1-\theta)]$（対数オッズ） \\
$t(y)$ & $y$ \\
$c(\phi)$ & $(1+e^\phi)^{-1}$ \\
$h(y)$ & $1$ \\
\bottomrule
\end{tabular}
\end{center}

\medskip
\noindent\textbf{共役事前分布の導出：}

$\phi$ に対する共役事前分布は
$p(\phi|n_0, t_0) \propto (1+e^\phi)^{-n_0} e^{n_0 t_0 \phi}$．

変数変換 $\phi \to \theta$ により $\theta$ についての事前分布を求めると：

\[
  p(\theta|n_0, t_0) \propto \theta^{n_0 t_0 - 1}(1-\theta)^{n_0(1-t_0)-1},
\]

これは $\mathrm{Beta}(n_0 t_0,\; n_0(1-t_0))$ 分布である．

\medskip
\noindent\textbf{解釈：}
\begin{itemize}
  \item $t_0$：$Y=1$ になる確率の事前期待値
  \item $n_0 = 1$（弱い事前情報）かつ $t_0 = 1/2$ なら
        $\mathrm{Beta}(1/2,\, 1/2)$（Jeffreys事前分布）
\end{itemize}

\medskip
\noindent\textbf{事後分布：}

$\mathrm{Beta}(t_0,\, 1-t_0)$ 事前分布（$n_0=1$）のもとで：

\[
  \left\{\theta\;\middle|\; y_1,\ldots,y_n\right\}
  \sim \mathrm{Beta}\!\left(t_0 + \sum y_i,\;\; (1-t_0) + \sum(1-y_i)\right).
\]

\end{exbox}

% ===========================
\section{例：ポアソンモデル}
% ===========================

\begin{exbox}{ポアソンモデル}

$Y \sim \mathrm{Poisson}(\theta)$ について：

\[
  p(y|\theta) = \frac{\theta^y e^{-\theta}}{y!}
  = \frac{1}{y!} \cdot e^{-e^\phi} \cdot e^{\phi y},
\]

ここで $\phi = \log\theta$．

\medskip
\noindent\textbf{各成分の対応：}

\medskip
\begin{center}
\begin{tabular}{ll}
\toprule
記号 & 内容 \\
\midrule
$\phi$ & $\log\theta$ \\
$t(y)$ & $y$ \\
$c(\phi)$ & $\exp(e^{-\phi}) = e^{-\theta}$ \\
$h(y)$ & $1/y!$ \\
\bottomrule
\end{tabular}
\end{center}

\medskip
\noindent\textbf{共役事前分布の導出：}

変数変換 $\phi \to \theta$ により $\theta$ の事前分布を求めると：

\[
  p(\theta|n_0, t_0) \propto \theta^{n_0 t_0 - 1} e^{-n_0 \theta},
\]

これは $\mathrm{Gamma}(n_0 t_0,\; n_0)$ 分布である．

\medskip
\noindent\textbf{解釈：}
\begin{itemize}
  \item $t_0$：$Y$ の事前期待値（母平均の事前予想）
  \item $n_0 = 1$（弱い事前情報）なら $\mathrm{Gamma}(t_0, 1)$
\end{itemize}

\medskip
\noindent\textbf{事後分布：}

$\mathrm{Gamma}(t_0, 1)$ 事前分布（$n_0=1$）のもとで：

\[
  \left\{\theta\;\middle|\; y_1,\ldots,y_n\right\}
  \sim \mathrm{Gamma}\!\left(t_0 + \sum y_i,\;\; 1 + n\right).
\]

\end{exbox}

% ===========================
\section{3モデルの対応まとめ}
% ===========================

\begin{center}
\renewcommand{\arraystretch}{1.8}
\begin{tabular}{lccc}
\toprule
& \textbf{二項モデル} & \textbf{ポアソンモデル} & \textbf{正規モデル（5.5節）} \\
\midrule
自然パラメータ $\phi$ & $\log\frac{\theta}{1-\theta}$ & $\log\theta$ & $(\theta/\sigma^2,\; -1/2\sigma^2)$ \\
十分統計量 $t(y)$ & $y$ & $y$ & $(y,\; y^2)$ \\
共役事前分布 & Beta & Gamma & Normal $\times$ Inv-Gamma \\
$t_0$ の解釈 & $E[Y]$ の事前予想 & $E[Y]$ の事前予想 & $(E[Y], E[Y^2])$ の事前予想 \\
$n_0$ の意味 & 事前サンプルサイズ & 事前サンプルサイズ & 事前サンプルサイズ \\
\bottomrule
\end{tabular}
\end{center}

% ===========================
\section{まとめ}
% ===========================

\begin{thmbox}{3.3節のまとめ}
\begin{enumerate}
  \item 二項・ポアソンモデルはどちらも\textbf{1次元指数型分布族}であり，
        密度が $p(y|\phi) = h(y)c(\phi)e^{\phi t(y)}$ の形で書ける．

  \item この構造があると，共役事前分布が\textbf{一般的な形}
        $p(\phi|n_0, t_0) \propto c(\phi)^{n_0} e^{n_0 t_0 \phi}$
        で構成でき，事後分布も同じクラスに収まる．

  \item 事後分布のパラメータは $n_0 \to n_0 + n$，
        $t_0 \to \dfrac{n_0 t_0 + n\bar{t}(y)}{n_0+n}$ という
        \textbf{単純な更新則}で得られる．

  \item ハイパーパラメータは $n_0$（事前サンプルサイズ），
        $t_0$（十分統計量の事前期待値）という\textbf{直感的な解釈}を持つ．

  \item この枠組みは5.5節で正規モデル（2次元指数型分布族）に拡張され，
        同じ発想で共役事前分布を期待値ベースで構成できることが示される．
\end{enumerate}
\end{thmbox}

\end{document}
